\section{Quantitative refinements and weighted profiles}\label{sec:extensions}
Throughout this section the coefficients form one nested Rademacher
sequence, except where another coefficient distribution is specified.
The estimates of the preceding sections allow the annular width to vary
with the degree, and their constants and thresholds are given in
Appendix~\ref{sec:constants}.

\subsection{A logarithmic almost-sure rate}\label{sec:log-rate}
We obtain the constant in Theorem~\ref{thm:log-rate} from the explicit
bounds below.
The deterministic analytic thresholds depend on
$C_*,c_{\rm HW}$, while the last exceptional block is random.
\begin{proof}[Proof of Theorem~\ref{thm:log-rate}]
The idea is to run the block argument of Section~\ref{sec:unit-circle} at
a width $K=K_N$ that grows with $N$, while $C_*$ and $c_{\rm HW}$ stay
fixed. By \eqref{eq:unit-unselected} the unselected roots contribute
$2\log2/K$ to the error in a block, so an error of order $1/\log n$
requires $K_N$ to grow at least like a multiple of $\log N$, while the
other constants grow exponentially in $K$, which bounds $K_N/\log N$ from
above. We derive these upper bounds, choose $K_N$, and then follow
Steps~1--5 of Section~\ref{sec:unit-circle} at the width $K_N$.

Table~\ref{tab:rate-constants} lists the constants of the argument that
depend on $K$, and $T=(\log N)/32$ is the clipping level of
Section~\ref{app:const-log}.
\begin{table}[ht]
\centering
\small
\begin{tabular}{@{}lll@{}}
\toprule
Quantity & Dependence on $K$ & Source\\
\midrule
$H_K$ & $10^8e^{20(K+1)}=H_0e^{20K}$ & Section~\ref{app:const-log}\\
$\sqrt{2+H_K}$ & $\le\sqrt{2+H_0}\,e^{10K}$ & $H_K$ and $2\le2e^{20K}$\\
$D_K$ & $10(K+1)^2e^{4K}$ & Section~\ref{app:const-profile}\\
$C_K$ & $\le C_{\rm bad}(K+1)^4e^{13K}$ & Section~\ref{app:growing-annulus}\\
$S_K$ & $S_0e^K$ & Section~\ref{app:growing-annulus}\\
Mesh mean threshold, 1st term & number times $(K+1)e^{8K}$ & Section~\ref{app:growing-annulus}\\
Mesh mean threshold, 2nd term & number times $(K+1)e^{11K}$ & Section~\ref{app:growing-annulus}\\
$C_1(K+2)$ & $\le C_{10}(K+1)$ & Section~\ref{app:growing-annulus}\\
Tail amplitude $a$ & factor $e^{2(K+2)}=e^4e^{2K}$ & Step~1, Section~\ref{sec:unit-circle}\\
\bottomrule
\end{tabular}
\caption{The constants of the proof in Section~\ref{sec:log-rate} that
depend on the width $K$. Here $H_0=10^8e^{20}$, and $S_0$, $C_{10}$,
$C_{\rm bad}$ and the numbers in the mesh mean threshold do not depend
on $K$. The constant $C_{\rm bad}$ depends only on $c_{\rm HW}$.}
\label{tab:rate-constants}
\end{table}
It follows that each bound in Table~\ref{tab:rate-width} is at most of
the form $C(K+1)^p(\log N)^qe^{bK}N^{-\beta}$, with $C$ independent of
$K$, real $p,q$, and $b,\beta>0$. Since
$e^{bK}N^{-\beta}=N^{bK/\log N-\beta}$, and since powers of
$K+1\le\log N$ and of $\log N$ are dominated by any power of $N$, if
$K\le c\log N$ then
\[
 C(K+1)^p(\log N)^qe^{bK}N^{-\beta}\quad
 \begin{cases}
  \text{tends to zero} & \text{if } c<\beta/b,\\
  \text{is summable along } N=j^8 & \text{if } c<(\beta-1/8)/b,
 \end{cases}
\]
where in the second case the exponent of $j=N^{1/8}$ stays below $-1$.
\begin{table}[ht]
\centering
\begin{tabular}{@{}llll@{}}
\toprule
Bound & $e^{bK}$ & $N^{-\beta}$ & $K/\log N$ below\\
\midrule
\multicolumn{4}{@{}l}{\emph{Must tend to zero}}\\
Lower logarithmic tail in $\varepsilon_K^*(N)$ & $e^{10K}$ & $N^{-1/32}$ & $1/320$\\
Mesh mean threshold, first term & $e^{8K}$ & $N^{-1/32}$ & $1/256$\\
Taylor ratio $4M_2a/d_0^2$ & $e^{3K}$ & $N^{-1/32}$ & $1/96$\\
$H_K^4/N$ & $e^{80K}$ & $N^{-1}$ & $1/80$\\
$N^{65/64}\cdot2a/d_0$ & $e^{2K}$ & $N^{-1/32}$ & $1/64$\\
Truncation bias $2H_KTN^{-1/2}$ in $\varepsilon_K^*(N)$ & $e^{20K}$ & $N^{-1/2}$ & $1/40$\\
Mesh mean threshold, second term & $e^{11K}$ & $N^{-7/16}$ & $7/176$\\
$(6400e^{8K})^2/N$ & $e^{16K}$ & $N^{-1}$ & $1/16$\\
$D_K/N$ and $2(1+2K)e^{4K}/N$ & $e^{4K}$ & $N^{-1}$ & $1/4$\\
\midrule
\multicolumn{4}{@{}l}{\emph{Must be summable along $N=j^8$}}\\
Occupation failure $H_KN^{-3/8}$ & $e^{20K}$ & $N^{-3/8}$ & $1/80$\\
Clipped-integral failure $\frac{H_K}{256}N^{-7/16}(\log N)^{-12}$ & $e^{20K}$ & $N^{-7/16}$ & $1/64$\\
Small-derivative failure $C_KN^{-3/8}(\log N)^4$ & $e^{13K}$ & $N^{-3/8}$ & $1/52$\\
\bottomrule
\end{tabular}
\caption{The bounds of the proof in Section~\ref{sec:log-rate} that
grow exponentially with the width $K$, and the upper bound on
$K/\log N$ that each requires: $\beta/b$ in the first part and
$(\beta-1/8)/b$ in the second.}
\label{tab:rate-width}
\end{table}
In the first part of Table~\ref{tab:rate-width}, the two terms of
$\varepsilon_K^*(N)$ that grow with $K$, and $D_K/N$, enter the bound on
$U_N/N$ below through $(8/K)(\varepsilon_K^*(N)+D_K/N)$, a term that is
$o(1/\log N)$ when $K/\log N$ has a positive limit and
$\varepsilon_K^*(N)+D_K/N$ tends to zero. Its other rows are ratios of
Table~\ref{tab:growing-annulus}, which must be at most one for the degree
thresholds of Section~\ref{app:growing-annulus} to hold. The second part
contains the exceptional probabilities in \eqref{eq:block-failure} that
depend on $K$. The remaining ratios of Table~\ref{tab:growing-annulus},
$2K/N$ and the excluded-sector ratio $80C_{10}(K+1)(\log N)N^{-15/32}$,
have no exponential factor and tend to zero whenever $K\le\log N$, and
its condition $\varepsilon_K^*(N)+D_K/N\le K/16$ holds for large $N$ once
$\varepsilon_K^*(N)$ and $D_K/N$ tend to zero.
The strongest condition is $K/\log N<1/320$, which comes from the lower
logarithmic tail. We choose
\[
 K_N=\left\lfloor\frac{\log N}{1000}\right\rfloor,
\]
where $N=j^8$ as in Section~\ref{sec:unit-circle}. Then
$K_N/\log N\le1/1000<1/320$, so that every condition in
Table~\ref{tab:rate-width} holds.
At $K=K_N$ the error bounds and deterministic thresholds of
Section~\ref{app:growing-annulus} become \eqref{eq:E2-majorants}, and
every threshold ratio tends to zero, so that the degree thresholds of
that section hold for all large $j$.
The block argument of Section~\ref{sec:unit-circle} applies at this width,
and its exceptional probability is bounded by \eqref{eq:block-failure}
with $K=K_N$ and $R=2$, that is, by
\[
 4p_{K_N}^*(N)+5p_2^*(N)+N^{-3}+2N^{-10}+C_{K_N}N^{-3/8}(\log N)^4,
\]
in the notation of Appendix~\ref{sec:constants}.
Here the term $N^{-3}$ is the local-count exception, with $K_N+2$ in place
of $K$. Since this parameter grows with $N$, we take it from
Lemma~\ref{lem:local-counts-general} with $a_k=1$, $B=1$, $\eta=1$ and
$\mathcal E=\Omega$. Then $A$ and $t_N=A\log N/N$ are as in the proof of
Lemma~\ref{lem:local-counts}, the thresholds $N\ge e$ and $t_N\le1$
do not depend on $K$, and the constant is $C_1(K_N+2)=(4A+4K_N+12)/\log2$.
Since $K_N\le(\log N)/1000$, we have
$e^{bK_N}\le e^{(b/1000)\log N}=N^{b/1000}$ for $b\ge0$, so that the
factors $e^{20K}$ of $H_K$ and $e^{13K}$ of $C_K$ give, with $N=j^8$,
\begin{align*}
 N^{-3/8}e^{20K_N}&\le N^{-3/8+20/1000}=N^{-71/200}=j^{-71/25},\\
 N^{-3/8}e^{13K_N}&\le N^{-3/8+13/1000}=N^{-181/500}=j^{-362/125}.
\end{align*}
These are the two slowly decaying powers in \eqref{eq:E2-majorants} that
come from the variable width.
By the third and fourth lines of \eqref{eq:E2-majorants}, with
$K_N+1\le\log N=8\log j$, the terms of \eqref{eq:block-failure} that
depend on $K_N$ satisfy
\begin{align*}
 &4p_{K_N}^*(N)+C_{K_N}N^{-3/8}(\log N)^4\\
 &\qquad\le4H_0N^{-71/200}+4N^{-12}
  +\frac{H_0}{64}N^{-167/400}(\log N)^{-12}\\
 &\qquad\quad+C_{\rm bad}(K_N+1)^4N^{-181/500}(\log N)^4\\
 &\qquad\le4H_0j^{-71/25}+4j^{-96}
  +\frac{H_0}{64}j^{-167/50}(8\log j)^{-12}\\
 &\qquad\quad+C_{\rm bad}(8\log j)^8j^{-362/125}.
\end{align*}
Their exponents in $j$ are less than $-1$, and their remaining factors are
polynomial in $\log j$.
The other terms, $5p_2^*(N)+N^{-3}+2N^{-10}$, do not depend on $K_N$, and
by Table~\ref{tab:probability-budget} at the width $2$ they are summable
along $N=j^8$.
Thus every term in \eqref{eq:block-failure} is summable.

The bounds
\eqref{eq:unit-block-length}--\eqref{eq:unit-parameters} of Step~1 in
Section~\ref{sec:unit-circle} are deterministic for $j\ge60$. At the
width $K_N$ the constant of Lemma~\ref{lem:second-derivative} is
$100e^{K_N+4}$, and \eqref{eq:unit-radius}, \eqref{eq:unit-parameters}
and $e^{bK_N}\le N^{b/1000}$ give
\begin{align*}
 \frac{2a}{d_0}&\le90e^{2(K_N+2)}N^{-67/64}\sqrt{\log N}
 \le90e^4N^{-67/64+1/500}\sqrt{\log N},\\
 N^{65/64}\,\frac{2a}{d_0}&\le90e^4N^{-1/32+1/500}\sqrt{\log N}
 =90e^4\sqrt{\log N}\,N^{-117/4000},\\
 \frac{4M_2a}{d_0^2}
 &\le180\cdot100e^{K_N+4}\,e^{2(K_N+2)}N^{-1/32}\log N\\
 &=18000e^8e^{3K_N}N^{-1/32}\log N
 \le18000e^8(\log N)N^{-113/4000}.
\end{align*}
Since the last two right-hand sides tend to zero, for all large $j$ we
have $2a/d_0<N^{-65/64}$, so that $2a/d_0\le1/N$, and $4M_2a\le d_0^2$.
These are the last two rows of Table~\ref{tab:growing-annulus}, in which
$180S_0e^4=18000e^8$.

Let
\[
 \delta_N=\max_{q=-1,0,1}
 \left|\frac{\nu_N(1+qN^{-65/64})}{N}-\frac12\right|.
\]
By the sparse estimate \eqref{eq:unit-thin-band}, whose five events are
taken at the fixed width $2$, we have
\[
 \delta_N\le2N^{-1/64}+C(2)N^{-1/64}(\log N)^7
 =O\bigl(N^{-1/64}(\log N)^7\bigr).
\]
Since $2a/d_0<N^{-65/64}$, the bound \eqref{eq:unit-crossings-centered}
on $c_N/N$ holds at the width $K_N$. Hence $|\nu_N(1)/N-1/2|\le\delta_N$
by the case $q=0$, and $c_N/N\le2\delta_N$ by
\eqref{eq:unit-crossings-centered} and the cases $q=\pm1$.

By the first line of \eqref{eq:unit-unselected}, the number $U_N$ of
unselected roots is bounded by the number of roots with
$||\alpha|-1|>K_N/N$ plus $B_{K_N,N}$. Outside the four radial events at
the width $K_N$, Lemma~\ref{lem:annular-mass} applies with
$\varepsilon=\varepsilon_{K_N}^*(N)$, by \eqref{eq:layer2-log}, and with
$D=D_{K_N}$, by Lemma~\ref{lem:profile-rate} and
Section~\ref{app:const-profile}, since
$\varepsilon_{K_N}^*(N)+D_{K_N}/N\le K_N/16\le K_N/8$ for large $N$.
Outside the event of \eqref{eq:layer2-bad} we have $B_{K_N,N}\le N^{31/32}$.
This is \eqref{eq:block-unselected} at $K=K_N$. Dividing it by $N$ and
using \eqref{eq:E2-majorants}, we get
\begin{align*}
 \frac{U_N}{N}
 &\le N^{-1/32}+\frac{2\log2}{K_N}
  +\frac8{K_N}\left(\varepsilon_{K_N}^*(N)+\frac{D_{K_N}}N\right)\\
 &\le N^{-1/32}+\frac{2\log2}{K_N}
  +\frac8{K_N}\left(E_0N^{-17/800}(\log N)^7
  +10(K_N+1)^2N^{-249/250}\right)\\
 &\le\frac{2\log2}{K_N}+o(1/\log N),
\end{align*}
since $K_N/\log N\to1/1000$, so that $\log N$ times each term other than
$2\log2/K_N$ tends to zero.
By the normalized matching inequality \eqref{eq:unit-block} of Step~4 in
Section~\ref{sec:unit-circle}, we therefore have, uniformly for
$N\le n\le N+m_j$,
\begin{align*}
 \left|\frac{\nu_n(1)}n-\frac12\right|
 &\le\left|\frac{\nu_N(1)}N-\frac12\right|
       +\frac{c_N+U_N}{N}+\frac{3m_j}{2N}
 \le3\delta_N+\frac{U_N}{N}+\frac{3m_j}{2N}\\
 &\le3\delta_N+\frac{2\log2}{K_N}+o(1/\log N)+\frac{27}2N^{-1/8}
 =\frac{2\log2}{K_N}+o(1/\log N),
\end{align*}
where $3m_j/(2N)\le\frac{27}2N^{-1/8}$ by \eqref{eq:unit-block-length},
and $(\log N)\delta_N\to0$.
With every error term written out, this inequality is
\eqref{eq:E2-final}.
Since the exceptional probabilities are summable, it follows from the
Borel--Cantelli lemma that, almost surely, only finitely many exceptional
events occur.
Since $m_j\le N$, we have $\log N\le\log n\le\log(2N)$, so that
$\log n/\log N\to1$ uniformly within the block. Since
$(\log N)/1000-1<K_N\le(\log N)/1000$, we also have
$K_N/\log N\to1/1000$ and $(\log N)\,2\log2/K_N\to2000\log2$.
Almost surely, every large $n$ lies in a block $N\le n\le N+m_j$ whose
exceptional event does not occur, and therefore, by the bound above in
the block of $n$,
\begin{align*}
 \limsup_{n\to\infty}(\log n)\left|\frac{\nu_n(1)}n-\frac12\right|
 &\le\limsup_{j\to\infty}\,\log(2N)
   \left(\frac{2\log2}{K_N}+o(1/\log N)\right)\\
 &=\limsup_{j\to\infty}\frac{\log(2N)}{\log N}
   \left((\log N)\frac{2\log2}{K_N}+o(1)\right)
 =2000\log2.
\end{align*}
\end{proof}

\subsection{Choice of the sparse degree sequence}
\begin{proposition}[Sparse-degree exponents]\label{prop:sparse-exponents}
With logarithmic cutoff $(\log N)/32$, deviation $N^{-1/32}$,
derivative parameter $L=N^{1/64}$, exceptional root count $N^{31/32}$,
and thin-band width $N^{-1-1/64}$, one may replace $j^8$ by
$\lfloor j^q\rfloor$ for every
\[
                         \frac83<q<16.
\]
In particular $j^3$ works. After retuning these powers, every real
$q>2$ works for the fixed-$K$ strong law and radial-profile proof.
For this construction, the second-moment conditions require $q>2$.
\end{proposition}
\begin{proof}
The idea is to rerun the block argument of Section~\ref{sec:unit-circle},
and its form at moving radii in Section~\ref{sec:moving-radii}, along the
degrees $\lfloor j^q\rfloor$, with the powers of the statement left as
parameters, and to collect the inequalities that each step imposes on
these powers and on $q$. We solve them for every $q>2$, and then for the
fixed powers of the statement.

In this proof we use $N=\lfloor j^q\rfloor$ and
$m=\lfloor(j+1)^q\rfloor-\lfloor j^q\rfloor$ in place of $j^8$ and $m_j$.
We claim that if $q>1$, then eventually $m\le(2+q2^q)N^{1-1/q}$ and
$m\le N$. Indeed, for every $j\ge1$,
\begin{align*}
 m&\le(j+1)^q-j^q+1\le q(j+1)^{q-1}+1\le q2^{q-1}j^{q-1}+1\\
 &\le q2^{q-1}(2N)^{1-1/q}+1\le q2^qN^{1-1/q}+1\le(2+q2^q)N^{1-1/q},
\end{align*}
where we use $\lfloor(j+1)^q\rfloor\le(j+1)^q$ and
$\lfloor j^q\rfloor>j^q-1$, the mean value theorem for the $q$-th power on
$[j,j+1]$, where its derivative is at most $q(j+1)^{q-1}$ since $q>1$,
and $j^{q-1}=(j^q)^{1-1/q}$ with $j^q<N+1\le2N$. The second bound
follows once $N^{1/q}\ge2+q2^q$. Also $m\ge1$, since
$(j+1)^q-j^q\ge qj^{q-1}\ge1$.
Since $\lfloor j^q\rfloor\ge j^q/2$, a bound on an exceptional
probability that is a power of $N$ times a power of $\log N$ is summable
over $j$ when $q$ times the exponent of $N$ is less than $-1$.

We choose a tail disk of radius $1+(K+2)/N$.
The maximal-tail estimate applies with
$a=15e^{2(K+2)}\sqrt{m\log N}$, and its exceptional event still has
probability at most $N^{-10}$. This is Lemma~\ref{lem:tails} with $K+2$
in place of $K$ and $m_{\max}=m$, which applies since $1\le m\le N$ and
$N\ge\max\{4,K+2\}$ for large $j$. By the bound on $m$, we have
\[
 a\le15e^{2(K+2)}(2+q2^q)^{1/2}N^{1/2-1/(2q)}\sqrt{\log N}.
\]

We take the derivative parameter $N^\ell$ in place of $L=N^{1/64}$, so
that $d_0=N^{3/2-\ell}$.
For the deterministic localization bounds we have
\begin{align*}
 \frac{2Na}{d_0}&=2N^{\ell-1/2}a=O_K(N^{\ell-1/(2q)}\sqrt{\log N}),\\
 \frac{4aN^{2\ell}}{N^3}\sup|f_N''|
    &\le\frac{4aN^{2\ell}}{N^3}\cdot100e^{K+4}N^{5/2}\sqrt{\log N}
    =O_K(N^{2\ell-1/(2q)}\log N),
\end{align*}
where the supremum is over $|z|\le1+(K+1)/N$ and is bounded by
Lemma~\ref{lem:second-derivative} outside its exceptional event, the two
$O_K$ bounds use the bound on $a$, and the implied constants also depend
on the fixed $q$.
It follows that Taylor isolation holds when $\ell<1/(4q)$. Indeed, then
$\ell-1/(2q)<2\ell-1/(2q)<0$, so both bounds tend to zero,
and for large $N$ they imply that $2a/d_0\le1/N$ and $4M_2a\le d_0^2$, where
$M_2=100e^{K+4}N^{5/2}\sqrt{\log N}$ and $4aN^{2\ell}/N^3=4a/d_0^2$.
A band of width $N^{-1-\kappa}$ contains the localization disks when
$\kappa<1/(2q)-\ell$, since
$N^{1+\kappa}\cdot2a/d_0=N^\kappa\cdot2Na/d_0
=O_K(N^{\kappa+\ell-1/(2q)}\sqrt{\log N})$.
The moving target in Theorem~\ref{thm:radial-profile} has scaled drift
$O_x(N^{-1/q})$, which the same band absorbs.
Indeed, for $N\le n\le N+m$ the target radius moves by
$|x/n-x/N|\le|x|m/N^2$, as in Section~\ref{sec:moving-radii}, and by the
bound on $m$ we have
\[
 N\cdot\frac{|x|m}{N^2}\le(2+q2^q)|x|N^{-1/q},\qquad
 N^{1+\kappa}\cdot\frac{|x|m}{N^2}\le(2+q2^q)|x|N^{\kappa-1/q}.
\]
Since $\kappa<1/(2q)-\ell<1/q$, the second right-hand side tends to
zero.

For the logarithmic clipping we set $T=t\log N$ and $u=N^{-w}$, and we
apply Remark~\ref{rem:T1-powers} with $t=w>0$, moment parameter
$6\log N$, and threshold $N^{-2t}$.
By \eqref{eq:point-bounds}, the one- and
two-point errors of Theorem~\ref{thm:T1} are at most $C(K)N^{-1/2}$, and
the removed angles and pairs there have measures at most $2N^{-1/2}$ and
$10N^{-1/2}$. The quantities $d_1$ and $v$ of the remark are the
one-point error plus the removed measure of angles, and the two-point
error plus the removed measure of pairs plus twice $d_1$. Both are
therefore at most $C(K)N^{-1/2}$, with a larger $C(K)$, and the remark
applies with its exponent equal to $1/2$ and with $C(K)$ for both of its
constants. Here the moment bound is \eqref{eq:layer2-nns}, whose constant
is $C_*$ and whose exponent is $6$, and the energy identity is
\eqref{eq:layer2-energy}. Therefore the logarithmic error satisfies
\begin{align*}
 \eta_{\log}
 &\le N^{-t}+2tC(K)(\log N)N^{-1/2}\\
 &\quad+e(12C_*\log N)^6\sqrt{2N^{-2t}+C(K)N^{-1/2}}+\tfrac32N^{-2t}
 =O_K(N^{-t}(\log N)^6)
\end{align*}
for $0<t\le1/4$, since then $N^{-1/2}\le N^{-2t}\le N^{-t}$, and
$1\le\log N\le(\log N)^6$ for $N\ge e$. The first condition of (S) below
implies that $t<1/8$.
The probabilities of the three exceptional events are bounded by
constants times $N^{-1/2+4t}$, $N^{-12}$ and $N^{-1/2+2t}(\log N)^2$, the
three terms of the probability bound of Remark~\ref{rem:T1-powers} with
$w=t$ and moment parameter $6\log N$.

The mesh geometry and the normal approximation apply to a general
derivative parameter $N^\ell$.
Indeed, in Section~\ref{app:const-small-derivatives} the bound
$D'_KNL^2\log N$ on the number of mesh points, the one-point probability
\eqref{eq:layer2-smallball} and the covariance bound
$(3\mathcal B_K+\mathcal T_K)N^{-1/2}$ for a pair are written with $L$
as a parameter, and the normal-approximation errors in them,
$\mathcal B_KN^{-1/2}$ and $\mathcal T_KN^{-1/2}$, do not involve $L$.
For a root-count threshold $N^{1-d}$, the mean and variance bounds in
Corollary~\ref{cor:T3-powers}, with both of its exponents equal to $1/2$
and its sector count $b_0=O_K(\sqrt N\log N)$, lead to conditions
(iii)--(vi) of the table below.
We use $d<1/2$ for the excluded-sector bound
$O_K(\sqrt N\log N)$, which counts at most $C(K+2)\log N$ roots in each
of $O(\sqrt N)$ disks by Lemma~\ref{lem:local-counts}, with $K+2$ in place of $K$.
For logarithmic concentration and the thin band, we impose
\[
 q\left(\frac12-4t\right)>1,\qquad
 q\left(\frac12-2t\right)>1,\qquad
 0<\kappa<\min\left(t,\frac1{2q}-\ell\right).
                                                        \tag{S}
\]
At the moving radii of Theorem~\ref{thm:radial-profile}, the probes of
Section~\ref{sec:moving-radii} have spacing $h=N^{-\kappa}$.
Fix an integer $R\ge|x|+3$, as in Section~\ref{sec:moving-radii}, let
$N\ge2R$, and let $y$ be one of the probes $x-h$, $x$ and $x+h$. We apply
Lemma~\ref{lem:local-radial-secant} at $y$, with $R$ in place of $K$,
which applies since $|y|+h\le|x|+2h<R<N$. We take
the centering $\frac12\log N$, the profile $F$ of
Proposition~\ref{prop:profiles}, $M_F=1/2$ and $B_F=1$, since
$0\le\Phi'\le1/2$ and $0\le\Phi\le1$ (Section~\ref{sec:moving-radii}). In
place of $\varepsilon$ we take the variance-profile error
$(R^2+e^{4R}/2)/N$ of Lemma~\ref{lem:profile-rate}, and we use
$\eta_{\log}=O_x(N^{-t}(\log N)^6)$, which is the bound on $\eta_{\log}$
above with $R$ in place of $K$, since $|y|<R$, so that its constant
depends on $x$ through $R$. Combined with
$|\Phi(y)-\Phi(x)|\le|y-x|/2$, the lemma yields
\begin{align*}
 \left|\frac{\nu_N(1+y/N)}N-\Phi(x)\right|
 &\le\frac{|y-x|}2+\frac h4+\frac RN
 +\frac{2(\eta_{\log}+(R^2+e^{4R}/2)/N)(1+R/N)}h\\
 &\le\frac74N^{-\kappa}
   +4N^\kappa\bigl(\eta_{\log}+(R^2+e^{4R}/2)N^{-1}\bigr)\\
 &=O_x\bigl(N^{-\kappa}+N^{\kappa-t}(\log N)^6\bigr),
\end{align*}
where we use $|y-x|\le h=N^{-\kappa}$, $R/N\le N^{-\kappa}$ and
$1+R/N\le2$ for large $N$.
This tends to zero since $\kappa<t$.
The strict power inequalities dominate the above logarithmic factors and
all fixed-$K$ constants.
We therefore obtain the nine conditions below, of which rows (ii), (vii),
(viii) and (ix) together form (S).
\begin{center}
\begin{tabular}{@{}lll@{}}
\toprule
 & Condition & Where it arises\\
\midrule
(i) & $\ell<1/(4q)$ & Taylor isolation\\
(ii) & $\kappa<1/(2q)-\ell$ & band containing the localization disks\\
(iii) & $0<d<4\ell$ & mean threshold of Corollary~\ref{cor:T3-powers}\\
(iv) & $d<\frac12-2\ell$ & mean threshold of Corollary~\ref{cor:T3-powers}\\
(v) & $d<\frac12$ & excluded sectors\\
(vi) & $q(\frac12-4\ell-2d)>1$ & root-count failure, summable\\
(vii) & $q(\frac12-4t)>1$ & occupation failure, summable\\
(viii) & $q(\frac12-2t)>1$ & clipped-integral failure, summable\\
(ix) & $0<\kappa<t$ & thin band, probes at the moving radii\\
\bottomrule
\end{tabular}
\end{center}

For every $q>2$ we choose
\begin{gather*}
 \Delta=\frac12-\frac1q,\qquad t=w=\frac\Delta{16},\\
 \ell=\min\left(\frac\Delta{32},\frac1{16q}\right),\qquad
 d=\ell,\qquad
 \kappa=\min\left(\frac t2,\frac1{8q}\right).
\end{gather*}
Then $0<\Delta<1/2$, all four powers are positive, and the nine
conditions follow from
\begin{alignat*}{3}
 &\text{(iii) }d=\ell<4\ell, &\qquad
 &\text{(v) }d\le\frac\Delta{32}<\frac12, &\qquad
 &\text{(iv) }d+2\ell=3\ell\le\frac{3\Delta}{32}<\frac12,\\
 &\text{(vii) }4t=\frac\Delta4<\Delta, &&
 \text{(viii) }2t=\frac\Delta8<\Delta, &&
 \text{(vi) }4\ell+2d=6\ell\le\frac{3\Delta}{16}<\Delta,\\
 &\text{(i) }2\ell\le\frac1{8q}<\frac1{2q}, &&
 \text{(ix) }\kappa\le\frac t2<t, &&
 \text{(ii) }\ell+\kappa\le\frac1{16q}+\frac1{8q}=\frac3{16q}<\frac1{2q},
\end{alignat*}
where the middle row implies the summability conditions (vii), (viii) and
(vi), since $q(\frac12-\Delta)=1$, and subtracting from $\frac12$ a
number smaller than $\Delta$ leaves more than $\frac12-\Delta$.
The remaining exceptional probabilities $N^{-3},N^{-10},N^{-12}$ are
summable as well. The conclusion then follows from the fixed-$K$
argument and the countable intersection over $K$.

For the fixed powers of the statement, $t=w=1/32$, $\ell=1/64$,
$d=1/32$ and $\kappa=1/64$, the conditions read
\begin{alignat*}{3}
 &\text{(vii) }q\left(\frac12-4t\right)=\frac{3q}8, &\qquad
 &\text{(vi) }q\left(\frac12-4\ell-2d\right)=\frac{3q}8, &\qquad
 &\text{(viii) }q\left(\frac12-2t\right)=\frac{7q}{16},\\
 &\text{(i) }\frac1{4q}-\ell=\frac{16-q}{64q}, &&
 \text{(ii) }\frac1{2q}-\ell-\kappa=\frac{16-q}{32q},
\end{alignat*}
where the three quantities of the first row must exceed one and the two
of the second row must be positive, while rows (ix), (iii) and (iv),
$\kappa<t$, $d<4\ell$ and $d<\frac12-2\ell$, hold since
$1/64<1/32<1/16<15/32$, as does (v).
For the fixed powers above, the slow exponent of the exceptional
probabilities is $3/8$, so that summability requires $q>8/3$. The
Taylor condition requires $q<16$.
At $q=8/3$ the mesh bound is of order $j^{-1}(\log j)^4$, and at $q=16$
the Taylor ratio grows logarithmically.
For $q\le2$ and $t>0$, the inequality $q(\frac12-4t)>1$ cannot hold,
since $q(\frac12-4t)<q/2\le1$.
These endpoint restrictions belong to the present second-moment estimates.
\end{proof}

\subsection{Sublevel components meeting nearby circles}\label{sec:nearby-circles}
\begin{corollary}[Sublevel components meeting a nearby circle]
\label{cor:crossings}
For a fixed tail bound $a_0$, and
$r_\pm=1\pm N^{-1-\kappa}$ with fixed $0<\kappa\le1/12$, the number
of roots in components of $\{|f_N|<a_0\}$ meeting those circles is
$o(N)$ outside a single exceptional sequence summable at $N=j^8$.
\end{corollary}
\begin{proof}
The idea is to bound, for each fixed $K$, the roots in question by the
roots of $f_N$ in a thin band about the unit circle, which
\eqref{eq:sign-thin-band} controls, and by the unselected roots $U_N$ of
Section~\ref{sec:unit-circle}, and then to let $K$ tend to infinity by
Lemma~\ref{lem:deterministic-diagonalization}. With $N=j^8$, let $Y_j$
denote the number of roots of $f_N$ in components of $\{|f_N|<a_0\}$ that
meet at least one of the circles $|z|=r_+$ and $|z|=r_-$. We show that
there exist deterministic numbers $\varepsilon_j\to0$ and events $F_j$ with
$\sum_j\Prob(F_j)<\infty$ such that $Y_j\le\varepsilon_jN$ off $F_j$.

Fix an integer $K\ge2$, and let $a=15e^{2(K+2)}\sqrt{m_j\log N}$ be the
maximal-tail level of Step~1 of Section~\ref{sec:unit-circle}.
Since $a\ge15\sqrt{\log N}$, we have $a_0\le a$ for all sufficiently
large $N$. Every component of $\{|f_N|<a_0\}$ lies in a component of
$\{|f_N|<a\}$, and if the first meets a circle, then so does the second.
Thus, for each circle, the number of roots in components of
$\{|f_N|<a_0\}$ meeting it is at most the number of roots in components
of $\{|f_N|<a\}$ meeting it.
A component containing a good annular root lies in the selected domain of
that root, which is a component of the sublevel set of Step~2 of
Section~\ref{sec:unit-circle}, whose level exceeds $a$. By that step, it
therefore has exactly one root and lies in the open disk of radius
$2a/d_0$ about it. Every other root is included in $U_N$.
If this component meets either nearby circle at a point $z$, then
$||\alpha|-1|\le|\alpha-z|+||z|-1|<N^{-1-\kappa}+2a/d_0$.
We enlarge the band to half-width $N^{-1-\kappa_*}$, which must contain
this range for large $N$ and to which \eqref{eq:sign-thin-band} must
apply. This requires, first, $0<\kappa_*<\kappa$, so that $N^{-1-\kappa}$
is small compared with $N^{-1-\kappa_*}$, second, $\kappa_*<3/64$, so
that $2a/d_0$ is small compared with $N^{-1-\kappa_*}$ by
\eqref{eq:unit-radius}, since $1-67/64=-3/64$, and third,
$\kappa_*\le1/64$, the range of \eqref{eq:sign-thin-band}.
We choose $\kappa_*=\min(\kappa/2,1/128)$. Since $\kappa_*\le\kappa/2$
and $\kappa_*\le1/128$, \eqref{eq:unit-radius} gives
\begin{align*}
 N^{1+\kappa_*}\bigl(N^{-1-\kappa}+2a/d_0\bigr)
 &\le N^{\kappa_*-\kappa}+90e^{2(K+2)}N^{1+\kappa_*-67/64}\sqrt{\log N}\\
 &\le N^{-\kappa/2}+90e^{2(K+2)}N^{-5/128}\sqrt{\log N},
\end{align*}
and the right-hand side tends to zero, so that
$N^{-1-\kappa}+2a/d_0<N^{-1-\kappa_*}$ for all $N$ large enough.
Since $0<\kappa_*\le1/128\le1/64$,
\eqref{eq:sign-thin-band} applies with $2$ in place of $K$ and $\kappa_*$
in place of $\kappa$. Outside five events, one at each radius
$1+qN^{-1-\kappa_*}$ with $q=-2,\dots,2$, it yields
\[
 \left|\frac{\nu_N(1\pm N^{-1-\kappa_*})}N-\frac12\right|
 \le2N^{-\kappa_*}+C(2)N^{\kappa_*-1/32}(\log N)^7.
\]
This five-radius estimate bounds the number of roots in this band by
$o(N)$:
\begin{align*}
 &\#\{\alpha:f_N(\alpha)=0,\ ||\alpha|-1|<N^{-1-\kappa_*}\}
 \le\nu_N(1+N^{-1-\kappa_*})-\nu_N(1-N^{-1-\kappa_*})\\
 &\qquad=N\left(\frac{\nu_N(1+N^{-1-\kappa_*})}N-\frac12\right)
   -N\left(\frac{\nu_N(1-N^{-1-\kappa_*})}N-\frac12\right)\\
 &\qquad\le2N\bigl(2N^{-\kappa_*}+C(2)N^{\kappa_*-1/32}(\log N)^7\bigr)
 =o(N),
\end{align*}
where the first inequality holds as in \eqref{eq:unit-crossings}, and the
last step since $\kappa_*>0$ and $\kappa_*-1/32\le1/128-4/128=-3/128$.
Each of the five events has probability at most the right-hand side of
\eqref{eq:sign-log-concentration} at $K=2$, which does not depend on the
radius. For fixed $K$, the exceptional probabilities are therefore
bounded by the right-hand side of \eqref{eq:unit-failure}, whose five
thin-band terms now bound the events at the exponent $\kappa_*$, and they
are summable. The explicit form of this bound is \eqref{eq:block-failure},
with thin-band parameter $\kappa_*$.
Let $E_{K,j}(\kappa_*)$ denote the event $E_{K,j}$ of Step~5 of
Section~\ref{sec:unit-circle} with its five thin-band events taken at the
exponent $\kappa_*$. Let $j_0(K)\ge60$ be a deterministic index beyond
which the thresholds of Section~\ref{sec:unit-circle} at the width $K$,
the bound $a_0\le a$, the strict inequality
$N^{-1-\kappa}+2a/d_0<N^{-1-\kappa_*}$ and the five-radius estimate hold.
Each root counted by $Y_j$ lies in a component of $\{|f_N|<a\}$ meeting
one of the circles. Hence it is either a good annular root, which lies in
the band $||\alpha|-1|<N^{-1-\kappa_*}$ by the strict inequality, or it
is counted by $U_N$. Therefore, for $j\ge j_0(K)$ and outside
$E_{K,j}(\kappa_*)$, by the band count and \eqref{eq:unit-unselected} we
have
\begin{align*}
 \frac{Y_j}N
 &\le\frac1N\#\{\alpha:f_N(\alpha)=0,\ ||\alpha|-1|<N^{-1-\kappa_*}\}
      +\frac{U_N}N\\
 &\le4N^{-\kappa_*}+2C(2)N^{\kappa_*-1/32}(\log N)^7\\
 &\quad+\frac{2\log2}{K}+C(K)N^{-1/32}(\log N)^7+N^{-1/32}.
\end{align*}

To apply Lemma~\ref{lem:deterministic-diagonalization}, which concerns
normalized block increments, with the exponent $8$, so that $n_j=j^8$,
we take $X_n=0$ when $n$ is an eighth power and $X_n=Y_j$ when
$j^8<n<(j+1)^8$. Since $X_{n_j}=0$ and each block contains more than one
degree, we have
$\max_{n_j\le n<n_{j+1}}|X_n-X_{n_j}|/n_j=Y_j/N$.
For each integer $K\ge2$ we take the index of the lemma to be $K-1$,
which runs over the positive integers, and put
\[
 b_{K-1}(j)=4N^{-\kappa_*}+2C(2)N^{\kappa_*-1/32}(\log N)^7
   +C(K)N^{-1/32}(\log N)^7+N^{-1/32},
\]
and let $p_{K-1}(j)$ be the bound \eqref{eq:unit-failure} for
$\Prob(E_{K,j}(\kappa_*))$ when $j\ge j_0(K)$, and $p_{K-1}(j)=1$ when
$j<j_0(K)$. By the last display,
$\Prob\{Y_j/N>2\log2/K+b_{K-1}(j)\}\le p_{K-1}(j)$ for every $j$.
Moreover $2\log2/K\to0$ as $K\to\infty$, and, for fixed
$K$, $b_{K-1}(j)\to0$ as $j\to\infty$, since $\kappa_*>0$ and
$\kappa_*-1/32\le-3/128$. Also $\sum_jp_{K-1}(j)<\infty$, since only finitely
many terms equal one and the others are summable by Step~5 of
Section~\ref{sec:unit-circle}. These are the hypotheses of
Lemma~\ref{lem:deterministic-diagonalization} at the index $K-1$, with
$2\log2/K$ as the term that tends to zero as the index grows, and the
lemma gives deterministic numbers $\varepsilon_j\to0$ and events $F_j$, the
events $E_j$ of the lemma, with $\sum_j\Prob(F_j)<\infty$ such that
$Y_j\le\varepsilon_jN$ off $F_j$. This proves the conclusion for both
signs simultaneously.
\end{proof}

\subsection{Polynomially varying deterministic coefficients}
Let $(\xi_k)_{k\geq0}$ be random variables on one probability space,
let $(c_k)$ be deterministic complex weights, and consider the prefixes
\[
 f_N(z)=\sum_{k=0}^N\xi_kc_kz^k.
\]
Note that every degree uses the same coefficient sequence. For independent
centered coefficients with unit variance, the normalization is the
square root of $\sum_{k=0}^N|c_k|^2r^{2k}$.

\begin{proposition}[Polynomially weighted logarithmic variance]
\label{prop:weighted-log-variance-profile}
Fix $\tau>-1/2$, put $c_k=k^\tau$ for $k\geq1$,
and choose a fixed finite $c_0\ne0$. For the normalization
\[
 \sigma_{N,\tau}(r)^2=|c_0|^2+\sum_{k=1}^N k^{2\tau}r^{2k},
\]
we have
\[
 \log\sigma_{N,\tau}(1+x/N)-(\tau+1/2)\log N
 \longrightarrow F_\tau(x):=\frac12\log\int_0^1t^{2\tau}e^{2xt}\dd t
\]
uniformly on compact intervals of $x$. Moreover,
\[
 F_\tau'(x)=
 \frac{\int_0^1t^{2\tau+1}e^{2xt}\dd t}
      {\int_0^1t^{2\tau}e^{2xt}\dd t},
 \qquad F_\tau'(0)=\frac{2\tau+1}{2\tau+2}.
\]
\end{proposition}

The four-secant formula in Lemma~\ref{thm:T2} applies, where
$F_{N,\tau}(x)=\log\sigma_{N,\tau}(1+x/N)-(\tau+1/2)\log N$.
With the centering $\frac12\log N$, we add $\tau\log N$ to $F_{N,\tau}$,
a term that cancels in every secant. Both the convention at $k=0$ and the
restriction $\tau>-1/2$ are needed for the integral profile.

\begin{proof}
Since every compact interval of $x$ lies in some $[-K,K]$, we fix $K>0$
and work with $|x|\leq K$. The idea is to show that
$\sigma_{N,\tau}(1+x/N)^2/N^{2\tau+1}$ converges to
$\int_0^1t^{2\tau}e^{2xt}\dd t$ uniformly for $|x|\leq K$, and then to
take logarithms. For $|x|\leq K$ and $0<t\leq1$, let
\[
 g_x(t)=t^{2\tau}e^{2xt},\qquad
 I_\tau(x)=\int_0^1g_x(t)\dd t=\int_0^1t^{2\tau}e^{2xt}\dd t.
\]
For the weighted sum, we have
\[
 S_{N,\tau}(x):=\frac1{N^{2\tau+1}}\sum_{k=1}^N k^{2\tau}e^{2xk/N}
 =\frac1N\sum_{k=1}^N(k/N)^{2\tau}e^{2xk/N}
 =\frac1N\sum_{k=1}^Ng_x(k/N).
\]

\emph{Step 1: the radial factor.}
Let $|x|\leq K$ and $N\geq2K$. Multiplying the comparison
\eqref{eq:profile-comparison} by $k^{2\tau}e^{2xk/N}/N^{2\tau+1}>0$ and
summing over $1\leq k\leq N$, we get, by the definition of $S_{N,\tau}$,
\[
 e^{-2K^2/N}S_{N,\tau}(x)
 \leq\frac1{N^{2\tau+1}}\sum_{k=1}^N k^{2\tau}(1+x/N)^{2k}
 \leq e^{2K^2/N}S_{N,\tau}(x).
\]

\emph{Step 2: the Riemann sums for $\tau\geq0$.}
When $\tau\geq0$, the integrand is continuous, since the map
$(x,t)\mapsto g_x(t)$ extends continuously to the compact set
$[-K,K]\times[0,1]$, and hence it is uniformly continuous there.
Let $\eta(\delta)$ be the supremum of $|g_x(t)-g_x(t')|$ over
$|x|\leq K$ and $t,t'\in[0,1]$ with $|t-t'|\leq\delta$, so that
$\eta(\delta)\to0$ as $\delta\to0$. For $|x|\leq K$ we have
\begin{align*}
 |S_{N,\tau}(x)-I_\tau(x)|
 &=\left|\sum_{k=1}^N\frac1Ng_x(k/N)
    -\sum_{k=1}^N\int_{(k-1)/N}^{k/N}g_x(t)\dd t\right|\\
 &=\left|\sum_{k=1}^N\int_{(k-1)/N}^{k/N}
    \bigl(g_x(k/N)-g_x(t)\bigr)\dd t\right|\\
 &\leq\sum_{k=1}^N\int_{(k-1)/N}^{k/N}|g_x(k/N)-g_x(t)|\dd t
 \leq\eta(1/N),
\end{align*}
where the last step uses $|k/N-t|\leq1/N$ on the $k$-th interval, whose
lengths add up to one.

\emph{Step 3: the Riemann sums for $-1/2<\tau<0$.}
When $-1/2<\tau<0$, we fix $0<\varepsilon<1$ and split the sum at
$k_\varepsilon=\lfloor\varepsilon N\rfloor$, where $N\geq2/\varepsilon$,
so that $2\leq k_\varepsilon<N$. Then
\begin{align*}
 S_{N,\tau}(x)-I_\tau(x)
 ={}&\left[\frac1N\sum_{k=1}^{k_\varepsilon}g_x(k/N)
      -\int_0^{k_\varepsilon/N}g_x(t)\dd t\right]\\
 &+\left[\frac1N\sum_{k=k_\varepsilon+1}^Ng_x(k/N)
      -\int_{k_\varepsilon/N}^1g_x(t)\dd t\right].
\end{align*}
Since $t\mapsto t^{2\tau}$ is decreasing, we have
$\frac1N(k/N)^{2\tau}\leq\int_{(k-1)/N}^{k/N}t^{2\tau}\dd t$ for
$k\geq2$, and the term $k=1$ is $\frac1N(1/N)^{2\tau}=N^{-(2\tau+1)}$.
Also $e^{2xt}\leq e^{2K}$ for $|x|\leq K$ and $0\leq t\leq1$, and
$k_\varepsilon/N\leq\varepsilon$. Hence the normalized initial part
satisfies
\begin{align*}
 0\leq\frac1N\sum_{k=1}^{k_\varepsilon}g_x(k/N)
 &\leq\frac{e^{2K}}N\sum_{k=1}^{k_\varepsilon}(k/N)^{2\tau}
 =e^{2K}\left(N^{-(2\tau+1)}
    +\sum_{k=2}^{k_\varepsilon}\frac1N(k/N)^{2\tau}\right)\\
 &\leq e^{2K}\left(N^{-(2\tau+1)}
    +\int_{1/N}^{k_\varepsilon/N}t^{2\tau}\dd t\right)
 \leq e^{2K}\left(\frac{\varepsilon^{2\tau+1}}{2\tau+1}
                  +N^{-(2\tau+1)}\right),
\end{align*}
where the intervals $[(k-1)/N,k/N]$ with $2\leq k\leq k_\varepsilon$ make
up $[1/N,k_\varepsilon/N]$.
By the same two facts, the corresponding integral satisfies
\[
 0\leq\int_0^{k_\varepsilon/N}g_x(t)\dd t
 \leq e^{2K}\int_0^\varepsilon t^{2\tau}\dd t
 =e^{2K}\frac{\varepsilon^{2\tau+1}}{2\tau+1}.
\]
Since both initial parts are nonnegative, the first bracket is at most
$e^{2K}(\varepsilon^{2\tau+1}/(2\tau+1)+N^{-(2\tau+1)})$ in absolute value.
For the remaining Riemann sum, every interval $[(k-1)/N,k/N]$ with
$k\geq k_\varepsilon+1$ lies in $[\varepsilon/2,1]$, since
$k_\varepsilon/N\geq\varepsilon-1/N\geq\varepsilon/2$. On the compact set
$[-K,K]\times[\varepsilon/2,1]$ the map $(x,t)\mapsto g_x(t)$ is
uniformly continuous. Let $\eta_\varepsilon(\delta)$ be the
supremum of $|g_x(t)-g_x(t')|$ over $|x|\leq K$ and
$t,t'\in[\varepsilon/2,1]$ with $|t-t'|\leq\delta$, so that
$\eta_\varepsilon(\delta)\to0$ as $\delta\to0$. Then the second bracket
satisfies
\begin{align*}
 &\left|\frac1N\sum_{k=k_\varepsilon+1}^Ng_x(k/N)
      -\int_{k_\varepsilon/N}^1g_x(t)\dd t\right|
 =\left|\sum_{k=k_\varepsilon+1}^N\int_{(k-1)/N}^{k/N}
    \bigl(g_x(k/N)-g_x(t)\bigr)\dd t\right|\\
 &\qquad\leq\sum_{k=k_\varepsilon+1}^N\int_{(k-1)/N}^{k/N}
    |g_x(k/N)-g_x(t)|\dd t
 \leq\eta_\varepsilon(1/N),
\end{align*}
where the last step holds as in Step~2, since $k/N$ and $t$ lie in
$[\varepsilon/2,1]$. Therefore
\[
 \sup_{|x|\leq K}|S_{N,\tau}(x)-I_\tau(x)|
 \leq e^{2K}\left(\frac{\varepsilon^{2\tau+1}}{2\tau+1}
                  +N^{-(2\tau+1)}\right)
     +\eta_\varepsilon(1/N).
\]
Since $2\tau+1>0$, the right-hand side tends to
$e^{2K}\varepsilon^{2\tau+1}/(2\tau+1)$ as $N\to\infty$, so that the limit
superior of the left side as $N\to\infty$ is at most this value for every
$0<\varepsilon<1$, and it is therefore zero.

We conclude that in both cases $S_{N,\tau}(x)\to I_\tau(x)$ uniformly for
$|x|\leq K$.

\emph{Step 4: the variance.}
Let $|x|\leq K$ and $N\geq2K$. By the definition of $\sigma_{N,\tau}$
and Step~1, we have
\begin{align*}
 \left|\frac{\sigma_{N,\tau}(1+x/N)^2}{N^{2\tau+1}}-I_\tau(x)\right|
 \leq{}&\frac{|c_0|^2}{N^{2\tau+1}}
   +\left|\frac1{N^{2\tau+1}}\sum_{k=1}^N k^{2\tau}(1+x/N)^{2k}
   -S_{N,\tau}(x)\right|\\
 &+|S_{N,\tau}(x)-I_\tau(x)|\\
 \leq{}&\frac{|c_0|^2}{N^{2\tau+1}}+(e^{2K^2/N}-1)S_{N,\tau}(x)
   +|S_{N,\tau}(x)-I_\tau(x)|\\
 \leq{}&\frac{|c_0|^2}{N^{2\tau+1}}
   +(e^{2K^2/N}-1)\left(\frac{e^{2K}}{2\tau+1}
   +|S_{N,\tau}(x)-I_\tau(x)|\right)\\
 &+|S_{N,\tau}(x)-I_\tau(x)|,
\end{align*}
where the second inequality also uses $1-e^{-2K^2/N}\leq e^{2K^2/N}-1$,
which holds since $e^u+e^{-u}\geq2$, and the last uses
$I_\tau(x)\leq e^{2K}\int_0^1t^{2\tau}\dd t=e^{2K}/(2\tau+1)$.
Since $2\tau+1>0$, the term with the fixed $c_0$ also vanishes after
normalization, and by Steps~2 and~3 the other terms tend to zero
uniformly for $|x|\leq K$, since $e^{2K^2/N}-1\to0$. Thus
$\sigma_{N,\tau}(1+x/N)^2/N^{2\tau+1}\to\int_0^1t^{2\tau}e^{2xt}\dd t$
uniformly for $|x|\leq K$.

\emph{Step 5: the logarithm.}
For $|x|\leq K$, the integral is at least $e^{-2K}/(2\tau+1)>0$, since
$I_\tau(x)\geq e^{-2K}\int_0^1t^{2\tau}\dd t$. By Step~1 we have
\begin{align*}
 \frac{\sigma_{N,\tau}(1+x/N)^2}{N^{2\tau+1}}
 &\geq\frac1{N^{2\tau+1}}\sum_{k=1}^Nk^{2\tau}(1+x/N)^{2k}
 \geq e^{-2K^2/N}S_{N,\tau}(x)\\
 &\geq e^{-2K^2/N}\bigl(I_\tau(x)-|S_{N,\tau}(x)-I_\tau(x)|\bigr)\\
 &\geq e^{-2K^2/N}\left(\frac{e^{-2K}}{2\tau+1}
      -|S_{N,\tau}(x)-I_\tau(x)|\right),
\end{align*}
where the first inequality drops $|c_0|^2/N^{2\tau+1}\geq0$.
Also $I_\tau(x)$ is at least the same lower bound, since
$e^{-2K^2/N}\leq1$. Since
$|\log a-\log b|\leq|a-b|/\min(a,b)$ for $a,b>0$, we get, for all $N$
so large that this lower bound is positive for $|x|\leq K$,
\begin{align*}
 &\left|\log\sigma_{N,\tau}(1+x/N)-(\tau+1/2)\log N-F_\tau(x)\right|\\
 &\qquad=\frac12\left|\log\frac{\sigma_{N,\tau}(1+x/N)^2}{N^{2\tau+1}}
    -\log I_\tau(x)\right|\\
 &\qquad\leq\frac{\bigl|\sigma_{N,\tau}(1+x/N)^2/N^{2\tau+1}-I_\tau(x)\bigr|}
   {2e^{-2K^2/N}\bigl(e^{-2K}/(2\tau+1)-|S_{N,\tau}(x)-I_\tau(x)|\bigr)},
\end{align*}
since $F_\tau=\frac12\log I_\tau$ and both arguments of the logarithms are
at least the lower bound above.
By Steps~2 to~4, the numerator tends to zero and the denominator tends
to $2e^{-2K}/(2\tau+1)$, both uniformly for $|x|\leq K$. Hence the left
side tends to zero uniformly for $|x|\le K$, which proves the profile
limit of the proposition.

By differentiating the positive integral, we get
\[
 \frac{\mathrm d}{\mathrm dx}\left(\frac12\log
                \int_0^1t^{2\tau}e^{2xt}\dd t\right)
 =\frac12\cdot\frac{\int_0^12t\cdot t^{2\tau}e^{2xt}\dd t}
        {\int_0^1t^{2\tau}e^{2xt}\dd t}
 =\frac{\int_0^1t^{2\tau+1}e^{2xt}\dd t}
        {\int_0^1t^{2\tau}e^{2xt}\dd t},
\]
where differentiation under the integral sign is allowed since, for
$|x|\leq K$, the $x$-derivative $2t^{2\tau+1}e^{2xt}$ of the integrand is
at most $2e^{2K}t^{2\tau}$, which is integrable on $[0,1]$ since
$2\tau>-1$. At zero this quotient is
$\int_0^1t^{2\tau+1}\dd t/\int_0^1t^{2\tau}\dd t=(2\tau+1)/(2\tau+2)$.
\end{proof}

\begin{remark}\label{rem:weighted-divergence}
For $\tau\leq-1/2$, the integral defining $F_\tau$ diverges at zero.
The finite-degree radial secant estimate still applies with the exact
normalization. To deduce an almost-sure root law from a variance profile,
one must also establish the occupation and logarithmic hypotheses and
the degree-block estimates.
\end{remark}
The integral diverges since its integrand is at least
$e^{-2|x|}t^{2\tau}$ and $2\tau\leq-1$. The finite-degree radial secant
estimate of the remark is \eqref{eq:T2-exact} of Lemma~\ref{thm:T2},
which allows any normalization $\sigma_N(r)>0$, and in particular
$\sigma_{N,\tau}$.



\begin{proposition}[Weighted polynomials with random signs]\label{prop:weighted-extension}
Take $c_0=1$ and $c_k=k^\tau$ for $k\ge1$, where $\tau>-1/2$.
For independent signs put
\[
 f_N^{(\tau)}(z)=\sum_{k=0}^N\epsilon_kc_kz^k,\qquad
 X_r(\theta)=\frac{f_N^{(\tau)}(re^{i\theta})}{\sigma_{N,\tau}(r)}.
\]
The uniform-on-compact variance limit in
Proposition~\ref{prop:weighted-log-variance-profile} holds for these polynomials.
Its candidate radial profile is
\[
 \Phi_\tau(x)=F_\tau'(x)
   =\frac{\int_0^1t^{2\tau+1}e^{2xt}\dd t}
          {\int_0^1t^{2\tau}e^{2xt}\dd t},\qquad
 \Phi_\tau(0)=\frac{2\tau+1}{2\tau+2}.
\]
The normalized logarithmic moments and angular energy satisfy
\[
 \mathbb E\int|\log|X_r(\theta)||^p\dd\mu(\theta)
       \le(C_*p)^{6p},\qquad p>1,
 \qquad
 \int|X_r(\theta)|^2\dd\mu(\theta)=1.
\]
The energy identity holds for every sign vector.
\end{proposition}
\begin{proof}
The idea is that $X_r$ is a Rademacher Fourier series whose deterministic
coefficients $c_kr^k/\sigma_{N,\tau}(r)$ have squared moduli summing to
one, which is all that the logarithmic bound and the energy identity of
Section~\ref{sec:log-integrability} use.
The variance convergence follows from the preceding proposition.
By the definition of $\sigma_{N,\tau}$ with $c_0=1$, the deterministic
Fourier coefficients satisfy
\[
 \sum_{k=0}^N\left|\frac{c_kr^k}{\sigma_{N,\tau}(r)}\right|^2
 =\frac{1+\sum_{k=1}^Nk^{2\tau}r^{2k}}{\sigma_{N,\tau}(r)^2}=1.
\]
By Corollary~1.2 of Nazarov--Nishry--Sodin, the above logarithmic
bound holds with the same absolute constant $C_*$.
Since that corollary holds for every $p\ge1$, the logarithmic bound also
holds for $p=1$, as in \eqref{eq:layer2-nns}.
By Parseval's identity, $\epsilon_k^2=1$ and the preceding identity, we
obtain for every sign vector the energy identity
\[
 \int|X_r(\theta)|^2\dd\mu(\theta)
 =\frac1{\sigma_{N,\tau}(r)^2}\sum_{k=0}^N\epsilon_k^2c_k^2r^{2k}=1.
\]
\end{proof}

The conclusions of Proposition~\ref{prop:weighted-extension} hold for any fixed finite nonzero $c_0$,
with $|c_0|^2$ in place of $1$ in the variance.
